\section{Suplemento sobre la cohomología de haces}
\label{section:0.12}
\label{section:0.12-fr}

\subsection{Cohomología de haces de módulos sobre espacios anillados}
\label{subsection:0.12.1}
\label{subsection:0.12.1-fr}

\begin{env}[12.1.1]
\label{0.12.1.1}
\label{0.12.1.1-fr}

Sea $(X,\mathcal O_X)$ un espacio anillado.
Recordemos que, para todo $\mathcal O_X$-módulo $\mathcal F$, se define la
\emph{cohomología} $H^\bullet(X,\mathcal F)$, que es un \emph{funtor cohomológico universal}
(T, 2.2) de la categoría $\boldsymbol{C}(X)$ de
$\mathcal O_X$-módulos en la categoría de grupos abelianos; es el funtor derivado
del funtor exacto por la izquierda
$\mathcal F\to\Gamma(X,\mathcal F)$.
El funtor $H^\bullet(X,\mathcal F)$ es isomorfo
\oldpage[0\textsubscript{III}]{50}
a la restricción a la categoría $\boldsymbol{C}(X)$ del funtor cohomológico
definido del mismo modo sobre la categoría de haces de grupos abelianos
en $X$ (G, II, 7.2.1).
\end{env}

\begin{env}[12.1.2]
\label{0.12.1.2}
\label{0.12.1.2-fr}
Pongamos $A=\Gamma(X,\mathcal O_X)$.
Como todo elemento de $A$ define un endomorfismo del grupo abeliano
$\Gamma(X,\mathcal F)$, define por funtorialidad un endomorfismo del
$\partial$-funtor $H^\bullet(X,\mathcal F)$.
Estos endomorfismos dotan a cada $H^p(X,\mathcal F)$ de una
estructura de $A$-\emph{módulo}, y el operador de conexión $\partial$ es
$A$-lineal. Además, para dos enteros positivos cualesquiera $p$, $q$, y dos $\mathcal O_X$-módulos $\mathcal F$, $\mathcal G$, se tiene un homomorfismo de $A$-módulos, llamado \emph{producto cup},
\[
\label{0.12.1.2.1-fr}
\tag{12.1.2.1}
H^p(X,\mathcal F)\otimes_A H^q(X,\mathcal G)
\longrightarrow H^{p+q}(X,\mathcal F\otimes_{\mathcal O_X}\mathcal G)
\]
(G, II, 6.6).
Estos homomorfismos hacen de la suma directa $S$ de los
$H^p(X,\mathcal O_X)$, para $p\geq0$, una $A$-\emph{álgebra graduada
anticomutativa}, y de la suma directa de los $H^p(X,\mathcal F)$ un
$S$-\emph{módulo graduado}.

Para todo recubrimiento abierto $\mathfrak U$ de $X$, denotaremos siempre
por $C^\bullet(\mathfrak U,\mathcal F)$ (contrariamente a
(G, II, 5.1)) el complejo de cocadenas alternadas del nervio de
$\mathfrak U$ con valores en el sistema de coeficientes
$\Gamma(U_\sigma,\mathcal F)$.
Está claro que
$C^\bullet(\mathfrak U,\mathcal F)$ es un
$A$-\emph{módulo graduado}, de modo que los grupos de cohomología
$H^p(\mathfrak U,\mathcal F)$ de este complejo están dotados de
una estructura de $A$-\emph{módulo}.
Además, las aplicaciones canónicas
$H^p(\mathfrak U,\mathcal F)\to H^p(X,\mathcal F)$
(G, II, 5.4) son $A$-\emph{lineales}.
\end{env}

\begin{env}[12.1.3]
\label{0.12.1.3}
\label{0.12.1.3-fr}
Sea $(X',\mathcal O_{X'})$ un segundo espacio anillado, y sea
$f=(\psi,\theta)$ un morfismo de $X'$ en $X$.
Pongamos $A'=\Gamma(X',\mathcal O_{X'})$; el $\psi$-morfismo $\theta$
define canónicamente un homomorfismo de anillos $A\to A'$.
Sea $\mathcal F$ un $\mathcal O_X$-módulo y $\mathcal F'$ un
$\mathcal O_{X'}$-módulo.
Para todo \emph{$f$-morfismo} $u:\mathcal F\to\mathcal F'$
\sref[0\textsubscript{I}]{0.4.4.1}, veremos que se puede definir,
para todo $p\geq0$, un dihomomorfismo
\[
\label{0.12.1.3.1-fr}
\tag{12.1.3.1}
u_p:H^p(X,\mathcal F)\longrightarrow H^p(X',\mathcal F').
\]

En efecto, como $\psi^*$ es exacto en la categoría de haces de grupos
abelianos sobre $X$, el funtor
$\mathcal F\leadsto H^\bullet(X',\psi^*(\mathcal F))$
es un $\partial$-funtor sobre dicha categoría, y se sabe que existe un
homomorfismo canónico de $\partial$-funtores
\[
\label{0.12.1.3.2-fr}
\tag{12.1.3.2}
H^\bullet(X,\mathcal F)\longrightarrow H^\bullet(X',\psi^*(\mathcal F))
\]
determinado únicamente por la condición de reducirse al homomorfismo canónico $\Gamma(X,\mathcal F)\to\Gamma(X',\psi^*(\mathcal F))$
(T, 3.2.2).
Además, todo elemento de $A$ determina un endomorfismo $\mu$ de
$\Gamma(X,\mathcal F)$ y un endomorfismo $\mu'$ de
$\Gamma(X',\psi^*(\mathcal F))$ tal que el diagrama
\[
\label{0.12.1.3.3-fr}
\tag{12.1.3.3}
\xymatrix{
\Gamma(X,\mathcal F)\ar[r]\ar[d]_\mu
&\Gamma(X',\psi^*(\mathcal F))\ar[d]^{\mu'}\\
\Gamma(X,\mathcal F)\ar[r]
&\Gamma(X',\psi^*(\mathcal F))
}
\]
es conmutativo.
\oldpage[0\textsubscript{III}]{51}
Por la propiedad de unicidad para prolongar morfismos de funtores
cohomológicos universales (T, 2.2), se obtienen prolongaciones únicas de
$\mu$ y $\mu'$ a la cohomología que hacen conmutativos los diagramas
análogos a \sref{0.12.1.3.3}.
Así, \sref{0.12.1.3.2} es un homomorfismo de $A$-\emph{módulos}.

Observemos ahora que
$f^*(\mathcal F)=\psi^*(\mathcal F)
\otimes_{\psi^*(\mathcal O_X)}\mathcal O_{X'}$,
y, por tanto, que existe un dihomomorfismo canónico
$\psi^*(\mathcal F)\to f^*(\mathcal F)$
del $\psi^*(\mathcal O_X)$-módulo $\psi^*(\mathcal F)$ en el
$\mathcal O_{X'}$-módulo $f^*(\mathcal F)$.
Por funtorialidad se obtiene, pues, un dihomomorfismo funtorial
\[
\label{0.12.1.3.4-fr}
\tag{12.1.3.4}
H^p(X',\psi^*(\mathcal F))\longrightarrow H^p(X',f^*(\mathcal F))
\]
siendo $A$ y $A'$ los anillos correspondientes.
Componiendo este dihomomorfismo con \sref{0.12.1.3.2}, se obtiene un
dihomomorfismo canónico, \emph{funtorial} en $\mathcal F$,
\[
\label{0.12.1.3.5-fr}
\tag{12.1.3.5}
\theta_p:H^p(X,\mathcal F)\longrightarrow H^p(X',f^*(\mathcal F)).
\]

Finalmente, por funtorialidad, el homomorfismo
$u^\sharp:f^*(\mathcal F)\to\mathcal F'$ da un homomorfismo de
$A'$-módulos
$H^p(X',f^*(\mathcal F))\to H^p(X',\mathcal F')$,
que, compuesto con \sref{0.12.1.3.5}, da
\sref{0.12.1.3.1}.

Sea $f'=(\psi',\theta'):X''\to X'$ un segundo morfismo de espacios
anillados, y sea $f''=f\circ f'$ el morfismo compuesto.
Teniendo en cuenta la compatibilidad del funtor $\psi'^*$ con los
productos tensoriales \sref[0\textsubscript{I}]{0.4.3.3}, se comprueba
inmediatamente que el compuesto del dihomomorfismo
$H^p(X',f^*(\mathcal F))\to H^p(X'',f'^*(f^*(\mathcal F)))$
con \sref{0.12.1.3.5} es el dihomomorfismo correspondiente
$H^p(X,\mathcal F)\to H^p(X'',f''^*(\mathcal F))$.
\end{env}

\begin{env}[12.1.4]
\label{0.12.1.4-fr}
Una definición directa del homomorfismo (\hyperref[0.12.1.3.2-fr]{12.1.3.2}) puede obtenerse de la manera siguiente: se considera una resolución inyectiva $\mathcal L^\bullet=(\mathcal L^i)$ de $\mathcal F$ constituida por haces de
grupos abelianos sobre $X$.
Como el funtor $\psi^*$ es exacto,
$\psi^*(\mathcal L^\bullet)$ es una \emph{resolución} de $\psi^*(\mathcal F)$
constituida por haces sobre $X'$.
Si $\mathcal L'^\bullet=(\mathcal L'^i)$ es una resolución inyectiva de
$\psi^*(\mathcal F)$ en la categoría de haces de grupos abelianos sobre
$X'$, existe un morfismo
$\psi^*(\mathcal L^\bullet)\to\mathcal L'^\bullet$
de complejos de haces de grupos abelianos, compatible con los
aumentos (M, V, 1.1.\emph{a)}, y determinado salvo homotopía.
Se obtienen así homomorfismos
\[
\Gamma(X,\mathcal L^\bullet)\longrightarrow
\Gamma(X',\psi^*(\mathcal L^\bullet))\longrightarrow
\Gamma(X',\mathcal L'^\bullet)
\]
de complejos de grupos abelianos.
Su compuesto, al pasar a la cohomología, da un morfismo de
$\partial$-funtores
$H^\bullet(X,\mathcal F)\to H^\bullet(X',\psi^*(\mathcal F))$; como coincide con (\hyperref[0.12.1.3.2-fr]{12.1.3.2}) en grado 0, es idéntico a él (T, 2.2).

Consideremos ahora un \emph{recubrimiento abierto}
$\mathfrak U=(U_\alpha)$ de $X$, y sea
$\mathfrak U'=(f^{-1}(U_\alpha))$ su recubrimiento abierto imagen inversa
de $X'$, imagen inversa de $\mathfrak U$. Los homomorfismos canónicos $\Gamma(V,\mathcal F)\to\Gamma(f^{-1}(V),f^*(\mathcal F))$ para todo abierto $V$ de $X$ definen inmediatamente (cf. G, II, 5.1) un homomorfismo de complejos
$C^\bullet(\mathfrak U,\mathcal F)\to
C^\bullet(\mathfrak U',f^*(\mathcal F))$,
y, por tanto, homomorfismos canónicos
\[
\label{0.12.1.4.1-fr}
\tag{12.1.4.1}
\theta_p:H^p(\mathfrak U,\mathcal F)
\longrightarrow H^p(\mathfrak U',f^*(\mathcal F)).
\]

\oldpage[0\textsubscript{III}]{52}
Además, existen diagramas conmutativos
\[
\label{0.12.1.4.2-fr}
\tag{12.1.4.2}
\xymatrix{
H^p(\mathfrak U,\mathcal F)\ar[r]^{\theta_p}\ar[d]
&H^p(\mathfrak U',f^*(\mathcal F))\ar[d]\\
H^p(X,\mathcal F)\ar[r]_{\theta_p}
&H^p(X',f^*(\mathcal F))
}
\]
donde las flechas verticales son los homomorfismos canónicos de
(G, II, 5.2).
Para demostrar la conmutatividad de \sref{0.12.1.4.2}, consideremos el
complejo de haces de cocadenas alternadas de $\mathcal F$ relativo a
$\mathfrak U$,
$\mathcal C^\bullet(\mathfrak U,\mathcal F)$, para el cual
$\Gamma(X,\mathcal C^\bullet(\mathfrak U,\mathcal F))=
C^\bullet(\mathfrak U,\mathcal F)$
(G, II, 5.2).
Los homomorfismos canónicos
$\Gamma(V,\mathcal F)\to
\Gamma(f^{-1}(V),\psi^*(\mathcal F))$
definen entonces un $\psi$-morfismo
$\mathcal C^\bullet(\mathfrak U,\mathcal F)\to
\mathcal C^\bullet(\mathfrak U',\psi^*(\mathcal F))$,
y, con la notación precedente, existe un diagrama conmutativo
\[
\xymatrix{
\Gamma(X,\mathcal C^\bullet(\mathfrak U,\mathcal F))\ar[r]\ar[d]
&\Gamma(X',\mathcal C^\bullet(\mathfrak U',\psi^*(\mathcal F)))\ar[d]\\
\Gamma(X,\mathcal L^\bullet)\ar[r]
&\Gamma(X',\psi^*(\mathcal L^\bullet))\ar[r]
&\Gamma(X',\mathcal L'^\bullet)
}
\]
Al pasar a la cohomología se obtienen diagramas conmutativos
\[
\xymatrix{
H^p(\mathfrak U,\mathcal F)\ar[r]\ar[d]
&H^p(\mathfrak U',\psi^*(\mathcal F))\ar[d]\\
H^p(X,\mathcal F)\ar[r]
&H^p(X',\psi^*(\mathcal F))
}
\]
donde las flechas verticales son los homomorfismos canónicos de
(G, II, 5.2).
Solo queda combinar estos diagramas con los diagramas conmutativos
\[
\xymatrix{
H^p(\mathfrak U',\psi^*(\mathcal F))\ar[r]\ar[d]
&H^p(\mathfrak U',f^*(\mathcal F))\ar[d]\\
H^p(X',\psi^*(\mathcal F))\ar[r]
&H^p(X',f^*(\mathcal F))
}
\]
\oldpage[0\textsubscript{III}]{53}
que provienen del homomorfismo
$\psi^*(\mathcal F)\to f^*(\mathcal F)$ y de la funtorialidad de los
homomorfismos canónicos de (G, II, 5.2), para obtener la conmutatividad
de \sref{0.12.1.4.2}.

Se observará que si
$A=\Gamma(X,\mathcal O_X)$ y $A'=\Gamma(X',\mathcal O_{X'})$, entonces
(12.1.4.3) es un dihomomorfismo de módulos correspondiente a
los anillos $A$ y $A'$.
Existe para \sref{0.12.1.4.1} una propiedad de transitividad respecto
de un compuesto de dos morfismos, análoga a la de
\sref{0.12.1.3.5}.
Finalmente, en las definiciones precedentes, en lugar de una resolución
inyectiva $\mathcal L^\bullet$ de $\mathcal F$, se habría podido partir igualmente de una resolución tal que $H^p(X,\mathcal L^i)=0$ para todo $i$ y todo $p>0$ (G, II, 4.7.1).
\end{env}

\begin{env}[12.1.5]
\label{0.12.1.5-fr}
Sean $\mathcal F$, $\mathcal G$, $\mathcal K$ tres $\mathcal O_X$-módulos, y consideremos
un $\mathcal O_X$-homomorfismo
$u:\mathcal F\otimes_{\mathcal O_X}\mathcal G\to\mathcal K$,
que da, en cohomología, los homomorfismos
\[
\label{0.12.1.5.1-fr}
\tag{12.1.5.1}
H^p(X,\mathcal F)\otimes_A H^q(X,\mathcal G)
\longrightarrow H^{p+q}(X,\mathcal K)
\]
deducidos del producto cup \sref{0.12.1.2.1}.
Demostraremos que, bajo las hipótesis y con la notación de
\sref{0.12.1.3}, los diagramas
\[
\label{0.12.1.5.2-fr}
\tag{12.1.5.2}
\xymatrix@C=10pt{
H^p(X,\mathcal F)\otimes_A H^q(X,\mathcal G)\ar[r]\ar[d]
&H^{p+q}(X,\mathcal K)\ar[d]\\
H^p(X',f^*(\mathcal F))\otimes_{A'}H^q(X',f^*(\mathcal G))\ar[r]
&H^{p+q}(X',f^*(\mathcal K))
}
\]
son conmutativos, donde las flechas verticales provienen de los homomorfismos canónicos (\hyperref[0.12.1.3.5-fr]{12.1.3.5}). Para ello, recordemos que (\hyperref[0.12.1.5.1-fr]{12.1.5.1}) puede obtenerse partiendo de las resoluciones \emph{canónicas} (G, II, 4.3) $\mathcal L^\bullet$, $\mathcal M^\bullet$, $\mathcal N^\bullet$ de $\mathcal F$, $\mathcal G$, $\mathcal K$ respectivamente, que están formadas por $\mathcal O_X$-módulos, y de la aplicación lineal
$\mathcal L^\bullet\otimes_{\mathcal O_X}\mathcal M^\bullet
\to\mathcal N^\bullet$
de complejos de $\mathcal O_X$-módulos correspondiente a $u$.
Esta aplicación proporciona un homomorfismo de complejos de $A$-módulos
$\Gamma(X,\mathcal L^\bullet)\otimes_A\Gamma(X,\mathcal M^\bullet)
\to\Gamma(X,\mathcal N^\bullet)$,
y, por tanto, al pasar a la cohomología, homomorfismos
$H^p(\Gamma(X,\mathcal L^\bullet))\otimes_A
H^q(\Gamma(X,\mathcal M^\bullet))
\to H^{p+q}(\Gamma(X,\mathcal N^\bullet))$
(G, II, 6.6).
Existe claramente un diagrama conmutativo
\[
\label{0.12.1.5.3-fr}
\tag{12.1.5.3}
\xymatrix@C=8pt{
\Gamma(X,\mathcal L^\bullet)\otimes_A\Gamma(X,\mathcal M^\bullet)
\ar[r]\ar[d]
&\Gamma(X,\mathcal N^\bullet)\ar[d]\\
\Gamma(X',\psi^*(\mathcal L^\bullet))
\otimes_{\Gamma(X',\psi^*(\mathcal O_X))}
\Gamma(X',\psi^*(\mathcal M^\bullet))\ar[r]
&\Gamma(X',\psi^*(\mathcal N^\bullet))
}
\]
\oldpage[0\textsubscript{III}]{54}
Al pasar a la cohomología se obtienen diagramas conmutativos
\[
\label{0.12.1.5.4-fr}
\tag{12.1.5.4}
\xymatrix@C=7pt{
H^p(X,\mathcal F)\otimes_A H^p(X,\mathcal G)\ar[r]\ar[d]
&H^{p+q}(X,\mathcal K)\ar[d]\\
H^p(\Gamma(X',\psi^*(\mathcal L^\bullet)))
\otimes_{\Gamma(X',\psi^*(\mathcal O_X))}
H^q(\Gamma(X',\psi^*(\mathcal M^\bullet)))\ar[r]
&H^{p+q}(\Gamma(X',\psi^*(\mathcal N^\bullet))).
}
\]
Pero como
$\psi^*(\mathcal L^\bullet)$,
$\psi^*(\mathcal M^\bullet)$, y
$\psi^*(\mathcal N^\bullet)$
son resoluciones de
$\psi^*(\mathcal F)$,
$\psi^*(\mathcal G)$, y
$\psi^*(\mathcal K)$ respectivamente, existe un diagrama conmutativo
(G, II, 6.6.1)
\[
\label{0.12.1.5.5-fr}
\tag{12.1.5.5}
\xymatrix@C=7pt{
H^p(\Gamma(X',\psi^*(\mathcal L^\bullet)))
\otimes_{\Gamma(X',\psi^*(\mathcal O_X))}
H^q(\Gamma(X',\psi^*(\mathcal M^\bullet)))\ar[r]\ar[d]
&H^{p+q}(\Gamma(X',\psi^*(\mathcal N^\bullet)))\ar[d]\\
H^p(X',\psi^*(\mathcal F))
\otimes_{\Gamma(X',\psi^*(\mathcal O_X))}
H^q(X',\psi^*(\mathcal G))\ar[r]
&H^{p+q}(X',\psi^*(\mathcal K))
}
\]
Finalmente, por funtorialidad, existe un diagrama conmutativo
\[
\label{0.12.1.5.6-fr}
\tag{12.1.5.6}
\xymatrix@C=7pt{
H^p(X',\psi^*(\mathcal F))
\otimes_{\Gamma(X',\psi^*(\mathcal O_X))}
H^q(X',\psi^*(\mathcal G))\ar[r]\ar[d]
&H^{p+q}(X',\psi^*(\mathcal K)).\ar[d]\\
H^p(X',f^*(\mathcal F))\otimes_{A'}H^q(X',f^*(\mathcal G))\ar[r]
&H^{p+q}(X',f^*(\mathcal K))
}
\]
Combinando los tres diagramas \sref{0.12.1.5.4},
\sref{0.12.1.5.5} y \sref{0.12.1.5.6}, se obtiene el diagrama
conmutativo deseado \sref{0.12.1.5.2}.
\end{env}

\begin{remark}[12.1.6]
\label{0.12.1.6}
\label{0.12.1.6-fr}
\oldpage[0\textsubscript{III}]{55}
Con la notación de \sref{0.12.1.3}, supongamos que existe un
diagrama conmutativo
\[
\label{0.12.1.6.1-fr}
\tag{12.1.6.1}
\xymatrix{
0\ar[r]&\mathcal F\ar[r]^r\ar[d]_u
&\mathcal G\ar[r]^s\ar[d]_v
&\mathcal H\ar[r]\ar[d]_w&0\\
0\ar[r]&\mathcal F'\ar[r]_{r'}
&\mathcal G'\ar[r]_{s'}
&\mathcal H'\ar[r]&0
}
\]
donde $r$, $s$ son homomorfismos de $\mathcal O_X$-módulos, $r'$, $s'$ homomorfismos de $\mathcal O_{X'}$-módulos, $u$, $v$, $w$ $f$-morfismos, y las filas son exactas.
Se sigue que el diagrama
\[
\label{0.12.1.6.2-fr}
\tag{12.1.6.2}
\xymatrix@C=6pt{
\ldots\ar[r]&H^p(X,\mathcal F)\ar[r]\ar[d]_{u_p}
&H^p(X,\mathcal G)\ar[r]\ar[d]_{v_p}
&H^p(X,\mathcal H)\ar[r]^\partial\ar[d]_{w_p}
&H^{p+1}(X,\mathcal F)\ar[r]\ar[d]_{u_{p+1}}&\ldots\\
\ldots\ar[r]&H^p(X',\mathcal F')\ar[r]
&H^p(X',\mathcal G')\ar[r]
&H^p(X',\mathcal H')\ar[r]_\partial
&H^{p+1}(X',\mathcal F')\ar[r]&\ldots
}
\]
es conmutativo.
En efecto, \sref{0.12.1.6.1} se factoriza como
\[
\xymatrix@R=14pt{
0\ar[r]&\mathcal F\ar[r]\ar[d]
&\mathcal G\ar[r]\ar[d]
&\mathcal H\ar[r]\ar[d]&0\\
0\ar[r]&\psi^*(\mathcal F)\ar[r]\ar[d]
&\psi^*(\mathcal G)\ar[r]\ar[d]
&\psi^*(\mathcal H)\ar[r]\ar[d]&0\\
0\ar[r]&\mathcal F'\ar[r]
&\mathcal G'\ar[r]
&\mathcal H'\ar[r]&0
}
\]
donde la fila intermedia es exacta
\sref[0\textsubscript{I}]{0.3.7.2}.
Basta entonces utilizar el hecho de que \sref{0.12.1.3.2} es un
homomorfismo de $\partial$-funtores y que los
$H^p(X',\mathcal F')$ forman un $\partial$-funtor en $\mathcal F'$.
\end{remark}

\begin{env}[12.1.7]
\label{0.12.1.7}
\label{0.12.1.7-fr}
Con las hipótesis y la notación de \sref{0.12.1.3}, consideremos ahora el
caso
$\mathcal F=f_*(\mathcal F')=\psi_*(\mathcal F')$; vamos a ver que el dihomomorfismo definido en (\hyperref[0.12.1.3-fr]{12.1.3}) \[
\label{0.12.1.7.1-fr}
\tag{12.1.7.1}
H^p(X,f_*(\mathcal F'))\longrightarrow H^p(X',\mathcal F')
\]
puede obtenerse (salvo un automorfismo de $H^p(X',\mathcal F')$) como un
homomorfismo de borde de una sucesión espectral para el funtor compuesto
$\mathcal F'\leadsto\Gamma(X',\psi_*(\mathcal F'))$ (T, 2.4). La descripción del homomorfismo (\hyperref[0.12.1.7.1-fr]{12.1.7.1}) dada en (\hyperref[0.12.1.4-fr]{12.1.4}) muestra aquí que se puede obtener este homomorfismo de la manera siguiente: se consideran resoluciones inyectivas $\mathcal L^\bullet$ y
$\mathcal L'^\bullet$ de $\psi_*(\mathcal F')$ y $\mathcal F'$,
respectivamente, y después un homomorfismo de complejos
$v:\psi^*(\mathcal L^\bullet)\to\mathcal L'^\bullet$
``por encima'' del homomorfismo canónico
$\psi^*(\psi_*(\mathcal F'))\to\mathcal F'$;
\oldpage[0\textsubscript{III}]{56}
se observa después que se tiene $\Gamma(X',\mathcal L'^\bullet)=\Gamma(X,\psi_*(\mathcal L'^\bullet))$
y que el homomorfismo compuesto
\[
\Gamma(X,\mathcal L^\bullet)\longrightarrow
\Gamma(X',\psi^*(\mathcal L^\bullet))
\xrightarrow{\Gamma(v)}\Gamma(X',\mathcal L'^\bullet)
\]
no es otro que
\[
\label{0.12.1.7.2-fr}
\tag{12.1.7.2}
\Gamma(v^\flat):\Gamma(X,\mathcal L^\bullet)
\longrightarrow\Gamma(X,\psi_*(\mathcal L'^\bullet))
\]
\sref[0\textsubscript{I}]{0.3.7.1}; y \sref{0.12.1.7.1} se
obtiene pasando a la cohomología en \sref{0.12.1.7.2}.

Por otra parte, las sucesiones espectrales para el funtor compuesto
$\mathcal F'\leadsto\Gamma(X,\psi_*(\mathcal F'))$
se obtienen tomando una resolución inyectiva de Cartan--Eilenberg
$\mathcal M^{\bullet\bullet}=(\mathcal M^{ij})$
del complejo $\psi_*(\mathcal L'^\bullet)$ en la categoría de
haces de grupos abelianos sobre $X$.
Las sucesiones espectrales en cuestión son las del bicomplejo
$\Gamma(X,\mathcal M^{\bullet\bullet})$; son birregulares, pues
$\mathcal M^{ij}=0$ para $i<0$ o $j<0$.
Ahora bien, la primera sucesión espectral de este bicomplejo degenera, pues
los haces $\psi_*(\mathcal L'^i)$ son flácidos (G, II, 3.1.1), por tanto $H_{\mathrm{II}}^q(\Gamma(X,\mathcal M^{i,\bullet}))
=H^q(\psi_*(\mathcal L'^i))=0$ para $q>0$ (G, II, 4.4.3); se tienen, pues, homomorfismos de borde \emph{biyectivos} (\hyperref[0.11.1.6-fr]{11.1.6})
\[
\label{0.12.1.7.3-fr}
\tag{12.1.7.3}
'E_2^{i0}=H^i(H_{\mathrm{II}}^0(\Gamma(X,\mathcal M^{\bullet\bullet})))
\longrightarrow H^i(\Gamma(X,\mathcal M^{\bullet\bullet}))
\]
y se sabe por \sref{0.11.3.4} que este homomorfismo proviene, al pasar
a la cohomología, del aumento
\[
\label{0.12.1.7.4-fr}
\tag{12.1.7.4}
\Gamma(X,\psi_*(\mathcal L'^\bullet))
\longrightarrow\Gamma(X,\mathcal M^{\bullet\bullet})
\]
que a su vez proviene del aumento
$\eta:\psi_*(\mathcal L'^\bullet)\to\mathcal M^{\bullet\bullet}$.

Para la segunda sucesión espectral, por otra parte, existen homomorfismos
de borde
\[
\label{0.12.1.7.5-fr}
\tag{12.1.7.5}
''E_2^{i0}=H^i(H_{\mathrm I}^0(\Gamma(X,\mathcal M^{\bullet\bullet}))
\longrightarrow H^i(\Gamma(X,\mathcal M^{\bullet\bullet}))
\]
que provienen, por \sref{0.11.3.4}, al pasar a la cohomología, del
homomorfismo de complejos
$Z_{\mathrm I}^0(\Gamma(X,\mathcal M^{\bullet\bullet}))
\to\Gamma(X,\mathcal M^{\bullet\bullet})$.
Como $\psi_*$ es exacto por la izquierda, la sucesión
\[
0\longrightarrow\psi_*(\mathcal F')\longrightarrow
\psi_*(\mathcal L'^0)\longrightarrow\psi_*(\mathcal L'^1)
\]
es exacta; por definición de una resolución de Cartan-Eilenberg (\hyperref[0.11.4.2-fr]{11.4.2}), se puede, pues, tomar $B_{\mathrm I}^0(\mathcal M^{\bullet\bullet})=0$, $Z_{\mathrm I}^0(\mathcal M^{\bullet\bullet})=\mathcal L^\bullet$; como el diagrama
\[
\xymatrix@C=50pt@R=24pt{
\mathcal L^0\ar[r]^{i^0}&\mathcal M^{00}\\
\psi_*(\mathcal F')\ar[u]^\varepsilon
\ar[r]_{\varepsilon'}\ar[ur]^{\varepsilon''}
&\psi_*(\mathcal L'^0)\ar[u]_{\eta^0}
}
\]
es conmutativo, la inyección de complejos
$i:\mathcal L^\bullet\to\mathcal M^{\bullet\bullet}$ es compatible con los
aumentos $\varepsilon$ y $\varepsilon''$.
Existen, pues, dos homomorfismos de complejos \textbf{de} $\mathcal L^\bullet$
en $\mathcal M^{\bullet\bullet}$,
\[
\xymatrix@C=52pt@R=22pt{
\mathcal L^\bullet\ar[r]^i\ar[dr]_{v^\flat}
&\mathcal M^{\bullet\bullet}\\
&\psi_*(\mathcal L'^\bullet)\ar[u]_\eta
}
\]
\oldpage[0\textsubscript{III}]{57}
compatibles con los aumentos $\varepsilon$ y
$\varepsilon''$.
Como $\mathcal L^\bullet$ es una resolución inyectiva y
$\mathcal M^{\bullet\bullet}$ está constituido por haces inyectivos, se sigue
de (M, V, 1.1 \textit{a)}) que estos dos homomorfismos son homotópicos.
Lo mismo vale, por tanto, para los dos homomorfismos correspondientes
$\Gamma(X,\mathcal L^\bullet)\to\Gamma(X,\mathcal M^{\bullet\bullet})$, y al pasar a la cohomología, se obtiene por tanto el \emph{mismo} homomorfismo; en otros términos, se ha mostrado que el homomorfismo de borde (\hyperref[0.12.1.7.5-fr]{12.1.7.5}), que se escribe $H^p(X,\psi_*(\mathcal F'))\to H^p(\Gamma(X,\mathcal M^{\bullet\bullet}))$
es el compuesto de \sref{0.12.1.7.1} y
\sref{0.12.1.7.3}, estando este último escrito
$H^p(X',\mathcal F')\to H^p(\Gamma(X,\mathcal M^{\bullet\bullet}))$ y que se ha visto que es un \emph{isomorfismo}; de donde nuestra aserción.
\end{env}

\subsection{Imágenes directas superiores}
\label{subsection:0.12.2-fr}

\begin{env}[12.2.1]
\label{0.12.2.1-fr}
Sean $(X,\mathcal O_X)$, $(Y,\mathcal O_Y)$ dos espacios anillados,
$f=(\psi,\theta)$ un morfismo de $X$ en $Y$, que define el funtor
\emph{imagen directa}
$f_*:\boldsymbol{C}(X)\longrightarrow\boldsymbol{C}(Y)$,
idéntico además a la restricción a $\boldsymbol{C}(X)$ del funtor $\psi_*$ definido
en la categoría de los haces de grupos abelianos sobre $X$. Este último
funtor es aditivo y exacto por la izquierda, y como todo haz de grupos
abelianos sobre $X$ es isomorfo a un subhaz de un haz \emph{inyectivo} de
grupos abelianos, se definen los \emph{funtores derivados} derechos
$\mathcal F\leadsto R^p\psi_*(\mathcal F)$ del funtor $\psi_*$; los
$R^p\psi_*(\mathcal F)$ son haces de grupos abelianos sobre $Y$, y
los $R^p\psi_*$ forman un \emph{funtor cohomológico universal} (T, 2.3).

Además, el haz $R^p\psi_*(\mathcal F)$ es el haz asociado al
prehaz $V\leadsto H^p(f^{-1}(V),\mathcal F)$ (T, 3.7.2). Si ahora
se supone que $\mathcal F$ es un $\mathcal O_X$-\emph{Módulo},
$H^p(f^{-1}(V),\mathcal F)$ está naturalmente provisto de una estructura de
$\Gamma(f^{-1}(V),\mathcal O_X)$-módulo, luego de
$\Gamma(V,\psi_*(\mathcal O_X))$-módulo, y el dato del homomorfismo
$\theta:\mathcal O_Y\to\psi_*(\mathcal O_X)$ permite deducir una
estructura de $\Gamma(V,\mathcal O_Y)$-módulo. Para las estructuras así
definidas, está claro que la restricción de un abierto $V$ a un abierto
$V'\subset V$ define un dihomomorfismo, y esto permite por tanto definir sobre
cada uno de los $R^p\psi_*(\mathcal F)$ una estructura de $\mathcal O_Y$-\emph{Módulo};
es este $\mathcal O_Y$-Módulo el que denotaremos por $R^p f_*(\mathcal F)$,
quedando así $R^p f_*$ definido como un funtor aditivo de $\boldsymbol{C}(X)$ en
$\boldsymbol{C}(Y)$. Además, los $R^p f_*$ forman un \emph{$\partial$-funtor}, pues si
$0\longrightarrow\mathcal F'\longrightarrow\mathcal F
\longrightarrow\mathcal F''\longrightarrow0$
es una sucesión exacta de $\mathcal O_X$-Módulos, la descripción de los
$R^p\psi_*$ y de la estructura de $\mathcal O_Y$-Módulo sobre
$R^p\psi_*(\mathcal F)$ dada arriba muestra inmediatamente que
el homomorfismo
$\partial:R^p\psi_*(\mathcal F'')\longrightarrow
R^{p+1}\psi_*(\mathcal F')$
es en este caso un homomorfismo de $\mathcal O_Y$-\emph{Módulos}. Finalmente, los
$R^p f_*$ se identifican con los funtores \emph{derivados} derechos de $f_*$: en efecto,
todo $\mathcal O_X$-Módulo admite una resolución \emph{inyectiva} formada por
$\mathcal O_X$-Módulos, y como tal resolución está formada por haces
\emph{flácidos} de grupos abelianos (G, II, 7.1), puede servir para calcular los
$R^p\psi_*(\mathcal F)$, puesto que $R^n\psi_*(\mathcal G)=0$ para $n\geq1$ y
para todo haz \emph{flácido} $\mathcal G$ (T, 2.4.1, Observación 3, y cor. de la
prop. 3.3.2). Se concluye pues que los $R^p f_*$ forman un \emph{funtor
cohomológico universal de} $\boldsymbol{C}(X)$ en $\boldsymbol{C}(Y)$ (T, 2.3).
\end{env}

\begin{env}[12.2.2]
\label{0.12.2.2}
\label{0.12.2.2-fr}
Sean $\mathcal F$ y $\mathcal G$ dos $\mathcal O_X$-módulos.
Con la notación de \sref{0.12.2.1}, para todo subconjunto abierto $V$ de
$Y$ existe el homomorfismo de producto cup \sref{0.12.1.2.1}
\[
H^p(f^{-1}(V),\mathcal F)
\otimes_{\Gamma(f^{-1}(V),\mathcal O_X)}
H^q(f^{-1}(V),\mathcal G)
\longrightarrow
H^{p+q}(f^{-1}(V),\mathcal F\otimes_{\mathcal O_X}\mathcal G)
\]
\oldpage[0\textsubscript{III}]{58}
Se sigue inmediatamente de la definición del producto cup
(G, II, 6.6) que estos homomorfismos conmutan con la restricción de
$V$ a un subespacio abierto $V'$ de $V$. Por otra parte, se tiene un homomorfismo de anillos
\[
\Gamma(V,\mathcal O_Y)\longrightarrow\Gamma(V,\psi_*(\mathcal O_X))
=\Gamma(f^{-1}(V),\mathcal O_X)
\]
proveniente de $\theta$, de donde un homomorfismo canónico de productos tensoriales
\[
\begin{aligned}
&H^p(f^{-1}(V),\mathcal F)\otimes_{\Gamma(V,\mathcal O_Y)}
H^q(f^{-1}(V),\mathcal G)\\
&\qquad\longrightarrow
H^p(f^{-1}(V),\mathcal F)
\otimes_{\Gamma(f^{-1}(V),\mathcal O_X)}
H^q(f^{-1}(V),\mathcal G)
\end{aligned}
\]
que es asimismo compatible con la restricción de $V'$ a $V$.
Por composición se obtiene así un homomorfismo de
$\Gamma(V,\mathcal O_Y)$-módulos.  Para los haces asociados a los
prehaces en cuestión, define un homomorfismo canónico,
funtorial en $\mathcal F$ y $\mathcal G$,
\[
\label{0.12.2.2.1-fr}
\tag{12.2.2.1}
R^p f_*(\mathcal F)\otimes_{\mathcal O_Y}R^q f_*(\mathcal G)
\longrightarrow
R^{p+q}f_*(\mathcal F\otimes_{\mathcal O_X}\mathcal G).
\]
Para $p=q=0$, este homomorfismo se reduce a
\sref[0\textsubscript{I}]{0.4.2.2.1}.
\end{env}

\begin{proposition}[12.2.3]
\label{0.12.2.3}
\label{0.12.2.3-fr}
Para todo $\mathcal O_X$-módulo $\mathcal F$ y todo $\mathcal O_Y$-módulo
$\mathcal L$ localmente libre de rango finito, existen isomorfismos canónicos
funtoriales
\[
\label{0.12.2.3.1-fr}
\tag{12.2.3.1}
R^p f_*(\mathcal F)\otimes_{\mathcal O_Y}\mathcal L
\xrightarrow{\sim}
R^p f_*(\mathcal F\otimes_{\mathcal O_X}f^*(\mathcal L)).
\]
\end{proposition}

\begin{proof}
El homomorfismo \sref{0.12.2.3.1} se obtiene componiendo el
siguiente caso particular de \sref{0.12.2.2.1},
\[
\label{0.12.2.3.2-fr}
\tag{12.2.3.2}
R^p f_*(\mathcal F)\otimes_{\mathcal O_Y}f_*(f^*(\mathcal L))
\longrightarrow
R^p f_*(\mathcal F\otimes_{\mathcal O_X}f^*(\mathcal L))
\]
con el homomorfismo del primer miembro de
\sref{0.12.2.3.1} al de \sref{0.12.2.3.2} inducido por el
homomorfismo canónico \sref[0\textsubscript{I}]{0.4.4.3.2}.
Para comprobar que \sref{0.12.2.3.1} es un isomorfismo cuando $\mathcal L$ es
localmente libre, puede reducirse inmediatamente al caso
$\mathcal L=\mathcal O_Y$, pues la cuestión es local sobre $Y$ y los
funtores en cuestión son aditivos en $\mathcal L$.
En ese caso, en vista de la definición de \sref{0.12.2.2.1}, la
proposición se reduce a comprobar que el homomorfismo correspondiente
de prehaces es biyectivo; esto se sigue inmediatamente de
$f^*(\mathcal O_Y)=\mathcal O_X$.
\end{proof}

\begin{env}[12.2.4]
\label{0.12.2.4}
\label{0.12.2.4-fr}
Sea $(Z,\mathcal O_Z)$ un tercer espacio anillado, y sea
$g:Y\to Z$ un morfismo de espacios anillados.
Se sabe (G, II, 7.1 y 3.1.1) que, para todo
$\mathcal O_X$-módulo inyectivo $\mathcal G$, $f_*(\mathcal G)$ es un haz flácido de grupos abelianos, y por consiguiente (\hyperref[0.12.2.1-fr]{12.2.1}) se tiene $R^p g_*(f_*(\mathcal G))=0$ para todo $p>0$. De ello se sigue (T, 2.4.1) que la \emph{sucesión espectral de Leray} de los funtores compuestos es aplicable al funtor compuesto $g_*f_*$: hay una sucesión espectral \emph{birregular} cuyo término límite es el funtor $R^\bullet h_*$, donde $h=g\circ f$, y cuyo término $E_2$ es
\[
\label{0.12.2.4.1-fr}
\tag{12.2.4.1}
E_2^{pq}=R^p g_*(R^q f_*(\mathcal F)).
\]
\end{env}

\begin{env}[12.2.5]
\label{0.12.2.5}
\label{0.12.2.5-fr}
Bajo las condiciones de \sref{0.12.2.4}, definiremos directamente
homomorfismos canónicos de $\mathcal O_Z$-módulos
\[
\label{0.12.2.5.1-fr}
\tag{12.2.5.1}
R^n g_*(f_*(\mathcal F))\longrightarrow R^n h_*(\mathcal F)
\]
y
\[
\label{0.12.2.5.2-fr}
\tag{12.2.5.2}
R^n h_*(\mathcal F)\longrightarrow g_*(R^n f_*(\mathcal F))
\]
\oldpage[0\textsubscript{III}]{59}
que pueden identificarse con los ``homomorfismos de borde'' de la sucesión
espectral de Leray (cf.\ \sref{0.12.1.7}).
Basta trabajar con los prehaces a los que están asociados los haces de
imágenes directas superiores \sref{0.12.2.1}.
Para ello, sea $W$ un subconjunto abierto cualquiera de $Z$ y su imagen inversa $g^{-1}(W)$ en $Y$.  Existe un
dihomomorfismo canónico
\[
\label{0.12.2.5.3-fr}
\tag{12.2.5.3}
H^n(g^{-1}(W),f_*(\mathcal F))\longrightarrow
H^n(f^{-1}(g^{-1}(W)),f^*(f_*(\mathcal F)))
\]
siendo los anillos correspondientes
$\Gamma(g^{-1}(W),\mathcal O_Y)$ y
$\Gamma(h^{-1}(W),\mathcal O_X)$.
Por otra parte, el homomorfismo canónico
\sref[0\textsubscript{I}]{0.4.4.3.3} da, por funtorialidad,
homomorfismos canónicos
\[
\label{0.12.2.5.4-fr}
\tag{12.2.5.4}
H^n(h^{-1}(W),f^*(f_*(\mathcal F)))\longrightarrow
H^n(h^{-1}(W),\mathcal F)
\]
que son homomorfismos de
$\Gamma(h^{-1}(W),\mathcal O_X)$-módulos.
Teniendo en cuenta el homomorfismo de anillos
$\Gamma(W,\mathcal O_Z)\to\Gamma(g^{-1}(W),\mathcal O_Y)$, se ve que, componiendo (\hyperref[0.12.2.5.4-fr]{12.2.5.4}) y (\hyperref[0.12.2.5.3-fr]{12.2.5.3}), se obtiene un homomorfismo de prehaces, que proporciona el homomorfismo de haces (\hyperref[0.12.2.5.1-fr]{12.2.5.1}).

La definición de (\hyperref[0.12.2.5.2-fr]{12.2.5.2}) es aún más simple; por definición, $R^n h_*(\mathcal F)$ está asociado al prehaz
$W\leadsto H^n(f^{-1}(g^{-1}(W)),\mathcal F)$
mientras que $R^n f_*(\mathcal F)$ está asociado al prehaz
$V\leadsto H^n(f^{-1}(V),\mathcal F)$;
Existe, por tanto, un homomorfismo canónico
\[
H^n(f^{-1}(g^{-1}(W)),\mathcal F)\longrightarrow
\Gamma(g^{-1}(W),R^n f_*(\mathcal F)),
\]
\textbf{y} es inmediato que estos homomorfismos definen un homomorfismo de
prehaces, que a su vez define \sref{0.12.2.5.2}.
\end{env}

\begin{env}[12.2.6]
\label{0.12.2.6}
\label{0.12.2.6-fr}
Bajo las hipótesis de \sref{0.12.2.4}, sean $\mathcal F$, $\mathcal G$
y $\mathcal K$ tres $\mathcal O_X$-módulos, y sea
$u:\mathcal F\otimes\mathcal G\to\mathcal K$
un $\mathcal O_X$-homomorfismo.
Entonces los diagramas siguientes son conmutativos:
\[
\xymatrix@C=10pt@R=24pt{
R^p g_*(f_*(\mathcal F))\otimes_{\mathcal O_Z}R^q g_*(f_*(\mathcal G))
\ar[r]\ar[d]&R^{p+q}g_*(f_*(\mathcal K))\ar[d]\\
R^p h_*(\mathcal F)\otimes_{\mathcal O_Z}R^q h_*(\mathcal G)
\ar[r]&R^{p+q}h_*(\mathcal K)
}
\tag{12.2.6.1}\label{0.12.2.6.1-fr}
\]
\textbf{y}
\[
\xymatrix@C=10pt@R=24pt{
R^p h_*(\mathcal F)\otimes_{\mathcal O_Z}R^q h_*(\mathcal G)
\ar[r]\ar[d]&R^{p+q}h_*(\mathcal K)\ar[d]\\
g_*(R^p f_*(\mathcal F))\otimes_{\mathcal O_Z}g_*(R^q f_*(\mathcal G))
\ar[r]&g_*(R^{p+q}f_*(\mathcal K))
}
\tag{12.2.6.2}\label{0.12.2.6.2-fr}
\]
\oldpage[0\textsubscript{III}]{60}
Aquí las flechas horizontales provienen de \sref{0.12.2.2.1} (la última
combinada con \sref[0\textsubscript{I}]{0.4.2.2.1}), y las flechas
verticales de \sref{0.12.2.5.1} y \sref{0.12.2.5.2}, respectivamente.

En efecto, basta verificar la aserción para los homomorfismos de prehaces
correspondientes.  Volviendo a las definiciones de
\sref{0.12.2.2} y \sref{0.12.2.5}, se reduce inmediatamente, para
\sref{0.12.2.6.1}, a los diagramas conmutativos
\sref{0.12.1.5.2}; la verificación es aún más sencilla para
\sref{0.12.2.6.2}.
\end{env}

\subsection{Suplementos sobre los funtores Ext de haces}
\label{subsection:0.12.3}
\label{subsection:0.12.3-fr}

\begin{env}[12.3.1]
\label{0.12.3.1-fr}
Consideremos un espacio anillado $(X,\mathcal O_X)$; no volveremos sobre la definición y las principales propiedades de los bifuntores $\operatorname{Ext}_{\mathcal O_X}^p(X;\mathcal F,\mathcal G)$ de la categoría de los $\mathcal O_X$-módulos en la de los $\Gamma(X,\mathcal O_X)$-módulos, y $\mathcal{E}xt_{\mathcal O_X}^p(\mathcal F,\mathcal G)$ de la categoría de los $\mathcal O_X$-módulos en sí misma, ni sobre la sucesión espectral birregular $E(\mathcal F,\mathcal G)$ que los relaciona (T, 4.2 y G, II, 7.3).
\end{env}

\begin{env}[12.3.2]
\label{0.12.3.2}
\label{0.12.3.2-fr}
Al igual que en (M, XIV, 1), se define la noción de una
\emph{extensión} de un $\mathcal O_X$-módulo $\mathcal F$ por un
$\mathcal O_X$-módulo $\mathcal G$ y la ley de composición sobre las clases de
extensiones equivalentes: los argumentos utilizados para los módulos se
trasladan manifiestamente a una categoría abeliana arbitraria.
La segunda demostración de (M, XIV, 1.1), que solo utiliza la existencia de
inmersiones en objetos inyectivos, sigue siendo, por tanto, válida en la
categoría de $\mathcal O_X$-módulos.  Muestra que
$\operatorname{Ext}_{\mathcal O_X}^1(X;\mathcal F,\mathcal G)$
se identifica canónicamente con el \emph{grupo abeliano de clases de
extensiones de $\mathcal F$ por $\mathcal G$}.
\end{env}

\begin{proposition}[12.3.3]
\label{0.12.3.3}
\label{0.12.3.3-fr}
Sea $(X,\mathcal O_X)$ un espacio anillado cuyo haz de anillos
$\mathcal O_X$ es coherente.
Entonces, para todo par de $\mathcal O_X$-módulos coherentes $\mathcal F$ y
$\mathcal G$ y todo $p\geq0$,
$\mathcal{E}xt_{\mathcal O_X}^p(\mathcal F,\mathcal G)$
es un $\mathcal O_X$-Módulo coherente.
\end{proposition}

\begin{proof}
Observemos que los $\mathcal{E}xt_{\mathcal O_X}^p(\mathcal F,\mathcal G)$ forman un funtor cohomológico
contravariante en $\mathcal F$.
Como $\mathcal F$ es coherente, para todo $p$ y todo punto $x\in X$
existen una vecindad abierta $U$ de $X$ y una sucesión exacta
de $(\mathcal O_X|U)$-módulos
\[
0\longrightarrow\mathcal R\longrightarrow\mathcal L_{p-1}
\longrightarrow\ldots\longrightarrow\mathcal L_0
\longrightarrow\mathcal F|U\longrightarrow0
\]
donde cada $\mathcal L_i$ ($0\leq i\leq p\setmin 1$) es isomorfo a
$\mathcal O_X^{n_i}|U$, y $\mathcal R$ es coherente.
Esto se sigue por inducción sobre $p$ de
\sref[0\textsubscript{I}]{0.5.3.2} y
\sref[0\textsubscript{I}]{0.5.3.4}, pues $\mathcal O_X$ es coherente. Observemos ahora que, para $p\geq1$,
$\operatorname{Ext}_{\mathcal O_X|U}^p(\mathcal L|U,\mathcal G|U)=0$
para todo $\mathcal O_X$-módulo $\mathcal L$ tal que $\mathcal L|U$ sea
isomorfo a $\mathcal O_X^n|U$ (T, 4.2.3).
El argumento de (M, V, 7.2) se aplica, por tanto, al funtor cohomológico
contravariante
$\mathcal F\leadsto
\mathrm{Hom}_{\mathcal O_X|U}(\mathcal F|U,\mathcal G|U)$
y da una sucesión exacta
\[
\mathcal{H}om_{\mathcal O_X|U}(\mathcal L_{p-1},\mathcal G|U)
\longrightarrow
\mathcal{H}om_{\mathcal O_X|U}(\mathcal R,\mathcal G|U)
\longrightarrow
\mathcal{E}xt_{\mathcal O_X|U}^p(\mathcal F|U,\mathcal G|U)
\longrightarrow0
\] y como los dos primeros términos de esta sucesión son $(\mathcal O_X|U)$-módulos coherentes
\sref[0\textsubscript{I}]{0.5.3.5}; por tanto, también lo es el tercero
\sref[0\textsubscript{I}]{0.5.3.4}.
\end{proof}

\oldpage[0\textsubscript{III}]{61}
\begin{proposition}[12.3.4]
\label{0.12.3.4}
\label{0.12.3.4-fr}
Sea $f:X\to Y$ un morfismo \textup{plano} de espacios anillados, y sean
$\mathcal F$ y $\mathcal G$ dos $\mathcal O_Y$-módulos.
\begin{enumerate}
  \item[\rm(i)] Existe un homomorfismo de bifuntores cohomológicos
  \[
\label{0.12.3.4.1-fr}
\tag{12.3.4.1}
f^*(\mathcal{E}xt_{\mathcal O_Y}^{\bullet}(\mathcal F,\mathcal G))
\xrightarrow{\sim}
\mathcal{E}xt_{\mathcal O_X}^{\bullet}(f^*(\mathcal F),f^*(\mathcal G))
\]
  que en grado $0$ se reduce al homomorfismo canónico
  \sref[0\textsubscript{I}]{0.4.4.6}.

  \item[\rm(ii)] Existe un morfismo canónico de sucesiones espectrales
  \[
\label{0.12.3.4.2-fr}
\tag{12.3.4.2}
E(\mathcal F,\mathcal G)\longrightarrow
E(f^*(\mathcal F),f^*(\mathcal G))
\]
  cuya aplicación sobre los términos $E_2$ viene dada por los homomorfismos
  \[
\label{0.12.3.4.3-fr}
\tag{12.3.4.3}
H^p(Y,\mathcal{E}xt_{\mathcal O_Y}^q(\mathcal F,\mathcal G))
\longrightarrow
H^p(X,\mathcal{E}xt_{\mathcal O_X}^q(f^*(\mathcal F),f^*(\mathcal G)))
\]
  deducidos de \sref{0.12.3.4.1} y \sref{0.12.1.3.1}.
\end{enumerate}
\end{proposition}

\begin{proof}
\begin{enumerate}
  \item[\rm(i)]
  Como $f^*$ es exacto sobre la categoría de $\mathcal O_Y$-módulos (\hyperref[0.6.7.2-fr]{\textbf{0}\textsubscript{I}, 6.7.2}), los funtores $\mathcal G\leadsto f^*(\mathcal{H}om_{\mathcal O_Y}(\mathcal F,\mathcal G))$ y $\mathcal G\leadsto\mathcal{H}om_{\mathcal O_X}(f^*(\mathcal F),f^*(\mathcal G))$ son exactos por la izquierda; se deduce canónicamente de (\hyperref[0.4.4.6-fr]{\textbf{0}\textsubscript{I}, 4.4.6}) un homomorfismo de sus funtores derivados. Para calcular estos últimos, se toma una resolución inyectiva $\mathcal L^\bullet=(\mathcal L^i)$ de $\mathcal G$.  Se tienen entonces
  homomorfismos
  $\mathcal H^p(f^*(\mathcal{H}om_{\mathcal O_Y}(\mathcal F,\mathcal L^\bullet)))\longrightarrow\mathcal H^p(\mathcal{H}om_{\mathcal O_X}(f^*(\mathcal F),f^*(\mathcal L^\bullet)))$
  entre los haces de cohomología de complejos.
  Por la exactitud de $f^*$, se tiene $\mathcal H^p(f^*(\mathcal{H}om_{\mathcal O_Y}(\mathcal F,\mathcal L^\bullet)))=f^*(\mathcal H^p(\mathcal{H}om_{\mathcal O_Y}(\mathcal F,\mathcal L^\bullet)))=f^*(\mathcal{E}xt_{\mathcal O_Y}^p(\mathcal F,\mathcal G))$ por definición. Por otra parte, la exactitud de $f^*$ implica también que
  $f^*(\mathcal L^\bullet)$ es una \emph{resolución} de $f^*(\mathcal G)$.
  Si $\mathcal L'^\bullet=(\mathcal L'^{i})$ es una resolución inyectiva
  de $f^*(\mathcal G)$ en la categoría de $\mathcal O_X$-módulos, existe un
  homomorfismo de complejos
  $f^*(\mathcal L^\bullet)\to\mathcal L'^\bullet$,
  determinado salvo homotopía, y por tanto un homomorfismo bien definido en
  cohomología.  Su compuesto con el homomorfismo precedente es
  \sref{0.12.3.4.1}.

  \item[\rm(ii)]
  Con la notación precedente, existe un homomorfismo de complejos de
  $\mathcal O_X$-módulos
  $f^*(\mathcal{H}om_{\mathcal O_Y}(\mathcal F,\mathcal L^\bullet))\longrightarrow\mathcal{H}om_{\mathcal O_X}(f^*(\mathcal F),\mathcal L'^\bullet)$.
  Sea $\mathcal M^{\bullet\bullet}$ una resolución inyectiva de
  Cartan--Eilenberg de
  $\mathcal{H}om_{\mathcal O_Y}(\mathcal F,\mathcal L^\bullet)$ en la categoría de
  $\mathcal O_Y$-módulos.  Como $f^*$ es exacto,
  $f^*(\mathcal M^{\bullet\bullet})$ es una resolución de Cartan--Eilenberg de
  $f^*(\mathcal{H}om_{\mathcal O_Y}(\mathcal F,\mathcal L^\bullet))$.
  Si $\mathcal M'^{\bullet\bullet}$ es una resolución inyectiva de
  Cartan--Eilenberg de
  $\mathcal{H}om_{\mathcal O_X}(f^*(\mathcal F),\mathcal L'^\bullet))$, existe
  \sref{0.11.4.2} un homomorfismo, determinado salvo homotopía,
  $f^*(\mathcal M^{\bullet\bullet})\longrightarrow\mathcal M'^{\bullet\bullet}$
  compatible con el homomorfismo precedente; en otras palabras, existe
  un $f$-morfismo
  $\mathcal M^{\bullet\bullet}\to
\mathcal M'^{\bullet\bullet}$ de
  de bicomplejos de haces.
  Induce un dihomomorfismo
  $\Gamma(Y,\mathcal M^{\bullet\bullet})\to
\Gamma(X,\mathcal M'^{\bullet\bullet})$
  de bicomplejos de módulos, determinado salvo homotopía, y por tanto
  un morfismo bien definido de sucesiones espectrales
  \sref{0.11.3.2}.  Este es precisamente el morfismo deseado
  \sref{0.12.3.4.2}, y la descripción
  \sref{0.12.3.4.3} se sigue inmediatamente de las definiciones.
\end{enumerate}
\end{proof}

\begin{proposition}[12.3.5]
\label{0.12.3.5}
\label{0.12.3.5-fr}
Bajo las hipótesis de \sref{0.12.3.4}, supongamos además que el
haz de anillos $\mathcal O_Y$ es coherente.
Entonces, para todo $\mathcal O_Y$-módulo coherente $\mathcal F$, los
homomorfismos canónicos \sref{0.12.3.4.1} son biyectivos.
\end{proposition}

\begin{proof}
\oldpage[0\textsubscript{III}]{62}
Como la cuestión es local sobre $Y$, se puede suponer que existe una
sucesión exacta
$0\longrightarrow\mathcal R\longrightarrow\mathcal O_Y^n\longrightarrow\mathcal F\longrightarrow0$,
Entonces $\mathcal R$ es también un $\mathcal O_Y$-módulo coherente
\sref[0\textsubscript{I}]{0.5.3.4}.
Para demostrar que los homomorfismos
\[
f^*(\mathcal{E}xt_{\mathcal O_Y}^p(\mathcal F,\mathcal G))
\longrightarrow
\mathcal{E}xt_{\mathcal O_X}^p(f^*(\mathcal F),f^*(\mathcal G))
\]
son biyectivos, razonamos por inducción sobre $p$.
Para $p=0$, la aserción es
\sref[0\textsubscript{I}]{0.6.7.6.1}.
Existe un diagrama conmutativo con filas exactas
\[
\hbox{\tiny$
\xymatrix@C=2pt@R=24pt{
f^*(\mathcal{E}xt_{\mathcal O_Y}^{p-1}(\mathcal O_Y^n,\mathcal G))
\ar[r]\ar[d]&
f^*(\mathcal{E}xt_{\mathcal O_Y}^{p-1}(\mathcal R,\mathcal G))
\ar[r]^\partial\ar[d]&
f^*(\mathcal{E}xt_{\mathcal O_Y}^{p}(\mathcal F,\mathcal G))
\ar[r]\ar[d]&
f^*(\mathcal{E}xt_{\mathcal O_Y}^{p}(\mathcal O_Y^n,\mathcal G))
\ar[d]\\
\mathcal{E}xt_{\mathcal O_X}^{p-1}(\mathcal O_X^n,f^*(\mathcal G))
\ar[r]&
\mathcal{E}xt_{\mathcal O_X}^{p-1}(f^*(\mathcal R),f^*(\mathcal G))
\ar[r]^\partial&
\mathcal{E}xt_{\mathcal O_X}^{p}(f^*(\mathcal F),f^*(\mathcal G))
\ar[r]&
\mathcal{E}xt_{\mathcal O_X}^{p}(\mathcal O_X^n,f^*(\mathcal G))
}
$}
\] puesto que $f^*(\mathcal O_Y)=\mathcal O_X$; como $f^*$ es exacto, las dos filas son exactas. Además, se tiene $\mathcal{E}xt_{\mathcal O_Y}^p(\mathcal O_Y^n,\mathcal G)=0$ para todo $p>0$ y análogamente $\mathcal{E}xt_{\mathcal O_X}^p(\mathcal O_X^n,f^*(\mathcal G))=0$ para todo $p>0$ (T, 4.2.3). Habida cuenta de la hipótesis de inducción, las dos primeras flechas verticales del diagrama precedente son isomorfismos, y los términos de la derecha son 0, por tanto
$f^*(\mathcal{E}xt_{\mathcal O_Y}^p(\mathcal F,\mathcal G))\longrightarrow\mathcal{E}xt_{\mathcal O_X}^p(f^*(\mathcal F),f^*(\mathcal G))$
es un isomorfismo.
\end{proof}

\subsection{Hipercohomología del funtor imagen directa}
\label{subsection:0.12.4}
\label{subsection:0.12.4-fr}

\begin{env}[12.4.1]
\label{0.12.4.1}
\label{0.12.4.1-fr}
Sean $(X,\mathcal O_X)$ y $(Y,\mathcal O_Y)$ dos espacios anillados, y sea
$f:X\to Y$ un morfismo de espacios anillados.
Se puede formar la hipercohomología de $f_*$ respecto de un complejo
\emph{arbitrario} de $\mathcal O_X$-módulos
$\mathcal K^\bullet=(\mathcal K^i)_{i\in\mathbb Z}$
(\hyperref[0.11.4.4-fr]{11.4.4}), pues en la categoría abeliana de los $\mathcal O_Y$-Módulos los límites inductivos filtrantes existen y son exactos (T, 3.1.1). Los $\mathcal O_Y$-módulos de hipercohomología
$\mathbf{R}^p f_*(\mathcal K^\bullet)$ se denotarán también
$\mathscr H^p(f,\mathcal K^\bullet)$, o simplemente
$\mathscr H_f^p(\mathcal K^\bullet)$.
Recordemos que $\mathscr H^\bullet(f,\mathcal K^\bullet)$ es la cohomología del
bicomplejo de $\mathcal O_Y$-módulos $f_*(\mathcal L^{\bullet\bullet})$,
donde $\mathcal L^{\bullet\bullet}$ es una resolución inyectiva de
Cartan--Eilenberg de $\mathcal K^\bullet$ en la categoría de
$\mathcal O_X$-módulos; $\mathscr H^\bullet(f,\mathcal K^\bullet)$ es el término límite de dos sucesiones espectrales ${}'\mathcal E(f,\mathcal K^\bullet)$ y ${}''\mathcal E(f,\mathcal K^\bullet)$ cuyos términos $E_2$
vienen dados por
\begin{equation}
\label{0.12.4.1.1-fr}
\tag{12.4.1.1}
{}'\mathcal E_2^{pq}=\mathcal H^p(\mathcal H^q(f,\mathcal K^\bullet))
\end{equation}
y
\begin{equation}
\label{0.12.4.1.2-fr}
\tag{12.4.1.2}
{}''\mathcal E_2^{pq}=\mathcal H^p(f,\mathcal H^q(\mathcal K^\bullet))
\quad(=R^p f_*(\mathcal H^q(\mathcal K^\bullet)))
\end{equation}
En estas fórmulas se ha adoptado la notación general $T(A^\bullet)$ para el transformado de un complejo por un funtor (\hyperref[0.11.2.1-fr]{11.2.1}), y se escribe
$\mathcal H^p(f,\mathcal F)$ en lugar de $R^p f_*(\mathcal F)$ para un $\mathcal O_X$-módulo $\mathcal F$. Recordemos además que la sucesión ${}'\mathcal E(f,\mathcal K^\bullet)$ es siempre \emph{regular}; las dos sucesiones espectrales ${}'\mathcal E(f,\mathcal K^\bullet)$ y ${}''\mathcal E(f,\mathcal K^\bullet)$ son \emph{birregulares} cuando $\mathcal K^\bullet$ está
acotado inferiormente,
\oldpage[0\textsubscript{III}]{63}
o cuando existe un entero $m$ tal que todo $\mathcal O_X$-módulo
admite una resolución flácida de longitud $\leq m$ \sref{0.11.4.4}.
\end{env}

\begin{env}[12.4.2]
\label{0.12.4.2}
\label{0.12.4.2-fr}
Denotaremos análogamente por $\mathbf H^\bullet(X,\mathcal K^\bullet)$ la
hipercohomología del funtor $\Gamma$ respecto de un complejo
$\mathcal K^\bullet$ de $\mathcal O_X$-módulos.
Así, los $\mathbf H^p(X,\mathcal K^\bullet)$ son
$\Gamma(X,\mathcal O_X)$-módulos.
Además, $\mathbf H^\bullet(X,\mathcal K^\bullet)$ puede considerarse como un
caso particular de $\mathscr H^\bullet(f,\mathcal K^\bullet)$, donde $f$ es un
morfismo de $(X,\mathcal O_X)$ a un espacio anillado constituido por un punto
dotado del anillo $\Gamma(X,\mathcal O_X)$.

Para todo subconjunto abierto $V$ de $X$, escribiremos
$\mathbf H^\bullet(V,\mathcal K^\bullet)$ en lugar de
$\mathbf H^\bullet(V,\mathcal K^\bullet|V)$.
\end{env}

\begin{proposition}[12.4.3]
\label{0.12.4.3}
\label{0.12.4.3-fr}
Para todo entero $p\in\mathbb Z$, el $\mathcal O_Y$-módulo
$\mathscr H^p(f,\mathcal K^\bullet)$ es canónicamente isomorfo al haz
asociado al prehaz
$U\leadsto\mathbf H^p(f^{-1}(U),\mathcal K^\bullet)$
sobre $Y$.
\end{proposition}

\begin{proof}
Con la notación de \sref{0.12.4.1}, el haz de cohomología
$\mathcal H^p(f_*(\mathcal L^{\bullet\bullet}))$ está asociado al
prehaz
$U\leadsto H^p(\Gamma(U,f_*(\mathcal L^{\bullet\bullet})))=H^p(\Gamma(f^{-1}(U),\mathcal L^{\bullet\bullet}))$.
Pero está claro que
$\mathcal L^{\bullet\bullet}|f^{-1}(U)$ es una resolución inyectiva de
Cartan--Eilenberg de $\mathcal K^\bullet|f^{-1}(U)$ (T, 3.1.3).
Por consiguiente,
$H^p(\Gamma(f^{-1}(U),\mathcal L^{\bullet\bullet}))=\mathbf H^p(f^{-1}(U),\mathcal K^\bullet)$
por definición.
\end{proof}

\begin{proposition}[12.4.4]
\label{0.12.4.4}
\label{0.12.4.4-fr}
La hipercohomología $\mathscr H^\bullet(f,\mathcal K^\bullet)$ es un
funtor cohomológico de $\mathcal K^\bullet$ en cada uno de los casos
siguientes:
\begin{enumerate}
  \item[\rm{a)}] $\mathcal K^\bullet$ varía en la categoría de complejos
    acotados inferiormente;
  \item[\rm{b)}] existe un entero $m$ tal que todo
    $\mathcal O_X$-módulo admite una resolución flácida de longitud $\leq m$;
  \item[\rm{c)}] $X$ es un espacio noetheriano.
\end{enumerate}
\end{proposition}

\begin{proof}
Los casos~\emph{a)} y~\emph{b)} son casos particulares de \sref{0.11.5.4}.
El caso~\emph{c)}, por otra parte, se sigue de \sref{0.11.5.2}, pues en
este caso el funtor $f_*$ conmuta con los límites inductivos
(G, II, 3.10.1).
\end{proof}

\begin{env}[12.4.5]
\label{0.12.4.5-fr}
Consideremos ahora un recubrimiento abierto $\mathfrak U=(U_\alpha)$ un recubrimiento abierto de $X$.
Para todo complejo de \emph{prehaces}
$\mathcal K^\bullet=(\mathcal K^j)$ sobre $X$, el bicomplejo $C^\bullet(\mathfrak U,\mathcal K^\bullet)$ cuya
componente de bigrado $(i,j)$ es
$C^i(\mathfrak U,\mathcal K^j)$, el grupo de
$i$-cocadenas alternadas del nervio de $\mathfrak U$ con valores
en $\mathcal K^j$ (G, II, 5.1).
Llamaremos a la cohomología de este bicomplejo la
\emph{hipercohomología} del recubrimiento $\mathfrak U$ con
coeficientes en $\mathcal K^\bullet$, y la denotaremos por
$\mathbf H^\bullet(\mathfrak U,\mathcal K^\bullet)=H^\bullet(C^\bullet(\mathfrak U,\mathcal K^\bullet))$.
La sucesión espectral de Leray de un recubrimiento
(T, 3.8.1 y G, II, 5.9.1) se generaliza a la hipercohomología como sigue.
\end{env}

\begin{proposition}[12.4.6]
\label{0.12.4.6}
\label{0.12.4.6-fr}
Sea $\mathcal K^\bullet=(\mathcal K^i)$ un complejo de
$\mathcal O_X$-módulos.
Existe un funtor espectral regular de $\mathcal K^\bullet$ con
término límite $\mathbf H^\bullet(X,\mathcal K^\bullet)$ y término $E_2$
\begin{equation}
\label{0.12.4.6.1-fr}
\tag{12.4.6.1}
\mathbf{E}_2^{pq}=\mathbf H^p(\mathfrak U,h^q(\mathcal K^\bullet))
\end{equation}
donde $h^q(\mathcal K^\bullet)$ denota el complejo de prehaces
$V\leadsto H^q(V,\mathcal K^\bullet)$
sobre $X$.
Esta sucesión espectral es birregular si $\mathcal K^\bullet$ está acotado inferiormente.
\end{proposition}

\begin{proof}
Consideremos una resolución inyectiva de Cartan-Eilenberg $\mathcal L^{\bullet\bullet}$ una resolución inyectiva de
Cartan--Eilenberg de $\mathcal K^\bullet$, y el tricomplejo $C^\bullet(\mathfrak U,\mathcal L^{\bullet\bullet})
=(C^i(\mathfrak U,\mathcal L^{jk}))$
Consideremos primero este tricomplejo como un bicomplejo en los grados $i$ y
$j+k$.
Como $i$ solo toma valores $\geq0$, la segunda sucesión espectral de este bicomplejo es regular (\hyperref[0.11.3.3-fr]{11.3.3}) y degenerada, pues se tiene $\mathbf H^q(\mathfrak U,\mathcal L^{jk})=0$ para todo $q>0$, siendo los $\mathcal O_X$-módulos $\mathcal L^{jk}$ haces flácidos
\oldpage[0\textsubscript{III}]{64}
(G, II, 5.2.3).
Se sigue que existe un isomorfismo canónico
$H^n(C^\bullet(\mathfrak U,\mathcal L^{\bullet\bullet}))\xrightarrow{\sim}H^n(\Gamma(X,\mathcal L^{\bullet\bullet}))$
\sref{0.11.1.6}, en virtud de (G, II, 5.2.2), y por tanto, por la
definición \sref{0.12.4.2}, un isomorfismo
$H^n(C^\bullet(\mathfrak U,\mathcal L^{\bullet\bullet}))\xrightarrow{\sim}\mathbf H^n(X,\mathcal K^\bullet)$.

Consideremos ahora
$C^\bullet(U,L^{\bullet\bullet})$
como un bicomplejo en los grados $i+j$ y $k$.
Como $k$ solo toma valores $\geq0$, la primera sucesión espectral de este bicomplejo es siempre regular; es birregular si $\mathcal L^{jk}=0$ para $j<j_0$, es decir, cuando
$\mathcal K^\bullet$ está acotado inferiormente \sref{0.11.3.3}.
Esta es la sucesión espectral deseada.
En efecto, para todo $j$, $\mathcal L^{j,\bullet}$ es una resolución
inyectiva de $\mathcal K^j$; por consiguiente,
$H^q(C^i(\mathfrak U,\mathcal L^{j,\bullet}))$
es precisamente el complejo de cocadenas
$C^i(\mathfrak U,h^q(\mathcal K^j))$.
Esto demuestra la proposición.
\end{proof}

\begin{corollary}[12.4.7]
\label{0.12.4.7}
\label{0.12.4.7-fr}
Si, para todo símplex $\sigma$ del nervio de $\mathfrak U$ y
todo entero $i$, se tiene $H^q(U_\sigma,\mathcal K^i)=0$ para $q>0$, entonces se tiene un isomorfismo canónico
\begin{equation}
\label{0.12.4.7.1-fr}
\tag{12.4.7.1}
\mathbf H^\bullet(\mathfrak U,\mathcal K^\bullet)
\xrightarrow{\sim}\mathbf H^\bullet(X,\mathcal K^\bullet).
\end{equation}
\end{corollary}

\begin{proof}
La hipótesis implica que
$C^\bullet(\mathfrak U,h^q(\mathcal K^j))=0$ para $q>0$, y por tanto que
$E_2^{pq}=0$ para $q>0$.
Como la sucesión espectral \sref{0.12.4.6.1} es degenerada y
regular, la conclusión se sigue de la definición
\sref{0.12.4.5} de
$\mathbf H^\bullet(\mathfrak U,\mathcal K^\bullet)$ y de
\sref{0.11.1.6}.
\end{proof}

\begin{env}[12.4.8]
\label{0.12.4.8}
\label{0.12.4.8-fr}
Sea $(X',\mathcal O_{X'})$ un segundo espacio anillado, y sea
$f=(\psi,\theta)$ un morfismo de $X'$ en $X$.
Por el mismo método que en \sref{0.12.1.3} y \sref{0.12.1.4}, se
define un dihomomorfismo para la hipercohomología de un complejo
$\mathcal K^\bullet$ de $\mathcal O_X$-módulos:
\begin{equation}
\label{0.12.4.8.1-fr}
\tag{12.4.8.1}
\mathbf H^p(X,\mathcal K^\bullet)
\longrightarrow\mathbf H^p(X',f^*(\mathcal K^\bullet)).
\end{equation}
Partiendo de una resolución inyectiva de Cartan--Eilenberg
$\mathcal L^{\bullet\bullet}$ de $\mathcal K^\bullet$, y como $\psi^*$ es exacto, $\psi^*(\mathcal L^{\bullet\bullet})$ es una
resolución de Cartan--Eilenberg de $\psi^*(\mathcal K^\bullet)$ en la
categoría de $\psi^*(\mathcal O_X)$-módulos.
Existe, por tanto, un morfismo
$\psi^*(\mathcal L^{\bullet\bullet})\to
\mathcal L'^{\bullet\bullet}$,
donde $\mathcal L'^{\bullet\bullet}$ es una resolución inyectiva de
Cartan--Eilenberg de $\psi^*(\mathcal K^\bullet)$.
Induce un morfismo en cohomología
$\mathbf H^\bullet(X,\mathcal K^\bullet)\longrightarrow\mathbf H^\bullet(X',\psi^*(\mathcal K^\bullet))$
Componiéndolo con el morfismo inducido por funtorialidad a partir de
$\psi^*(\mathcal K^\bullet)\to f^*(\mathcal K^\bullet)$, se obtiene el morfismo (\hyperref[0.12.4.8.1-fr]{12.4.8.1}) buscado.

Partiendo de (\hyperref[0.12.4.8.1-fr]{12.4.8.1}) y de (\hyperref[0.12.4.3-fr]{12.4.3}), se puede entonces, razonando como en (\hyperref[0.12.2.5-fr]{12.2.5}), definir, para dos morfismos $f:X\to Y$, $g:Y\to Z$ de espacios anillados, homomorfismos para la hipercohomología de un complejo $\mathcal K^\bullet$ de $\mathcal O_X$-módulos:
\begin{equation}
\label{0.12.4.8.2-fr}
\tag{12.4.8.2}
\mathscr H^n(g,f_*(\mathcal K^\bullet))
\longrightarrow\mathscr H^n(h,\mathcal K^\bullet)
\end{equation}
y
\begin{equation}
\label{0.12.4.8.3-fr}
\tag{12.4.8.3}
\mathscr H^n(h,\mathcal K^\bullet)
\longrightarrow g_*(\mathscr H^n(f,\mathcal K^\bullet)).
\end{equation}
Dejamos al lector los detalles de las definiciones.
\end{env}
